\section{Semi-Structured Interview Guide for Expert Interviews with Facilitators}
\label{annex:interview-guide-facilitators}

\subsection*{Einstiegstext (sinngemäß)}

Vielen Dank, dass Sie sich die Zeit nehmen. Ich untersuche, wie zielorientierte Remote-Meetings
so moderiert werden, dass die Gruppe beim vereinbarten Ziel bleibt und zu einem
verbindlichen Ergebnis kommt, also zu einer Entscheidung im engeren Sinn.

Gemeint sind damit Entscheidungssituationen, in denen die Gruppe zu einer Festlegung darüber gelangen
muss, wie es weitergeht. Die Art der Entscheidung kann dabei sehr vielfältig sein. Typische Fälle sind
zum Beispiel:
\begin{itemize}
  \item Priorisierung: z.B. Feature-Festlegung, entscheidung welche Vorhaben / welche Strategie am
    vielversprechensten ist
  \item Commitment: z.B. das Setzen eines ideal Team-Ziels, oder für das Team optimale Arbeitsweisen
  \item Oder Sonstinges in diesem Bereichen
\end{itemize}
Entscheidend ist, dass aus dem Meeting eine gemeinsame Festlegung folgt, an der sich die
weitere Arbeit orientiert.

Diese und andere Interviews dienen als Grundlage, um Anfoderungen für ein KI-Gestützes Meeingsystem
abzuleiten, welches Meeting-Teilnehmenden dabei untersützt sich in convergenten Zone des Meetings an das ziel zu halten.

Das Gespräch ist ein semi-strukturiertes Interview. Das heißt, dass ich bei Bedarf Nachfragen stelle, um
detailiertere Informationen zu erhalten.
Es gibt dabei aber keine richtigen oder falschen Antworten; mich interessiert Ihre Einschätzung basierend auf
ihre berfulcihen Praxis.
Das Interview dauert etwa 30 Minuten, wird über Zoom geführt und aufgezeichnet um späteres transkribiert zu
ermöglichen. Das Interview wird in anonymiserte bzw. pseudonymisierter Form in Rahmen dieser Arbeit ausgewertet.

Die Aufnahme wird nicht an Dritte weitergegeben. Sind Sie mit der Aufzeichnung und dieser
Verarbeitung einverstanden?

\textit{Bitte um klare Bestätigung, welche mit Aufgezeichnet wird.}

\medskip

% ---------------------------------------------------------------------------------------------------------------------
\subsection*{Beruflicher Kontext}
\textit{Funktion: Sampling-Validierung (Ca. 3 Minuten)}

\begin{enumerate}[label=1.\arabic*, leftmargin=*]
  \item Bitte beschreiben Sie kurz Ihre Erfahrungen im Bereich Moderation bzw. Facilitating, insbesondere in
    Situation in denen die Gruppe das Ziel hatte eine Entscheidung, wie z.B. das Festlegung nächster Schritte,
    Priorisierung oder Einigung über z.B. gemeinsame Vereinbarungen (z.B. Gemeinsame Arbeitsweisen,
    Verhaltens regeln) zu treffen.
    \begin{itemize}[label=\textendash]
      \item \textit{Optional:} Seit wann sind Sie in diesem Bereich tätig, und welche Teams oder
        Organsationen begleiten sie typischerweise?
    \end{itemize}
\end{enumerate}

% ---------------------------------------------------------------------------------------------------------------------
\subsection*{Herausforderungen in zielorientierten Remote-Meetings (RQ1)}
\textit{Funktion: Verdichtung der in Kapitel~2 identifizierten Herausforderungen (Ca. 6 Minuten)}

\begin{enumerate}[label=2.\arabic*, leftmargin=*]
  \item Wenn Sie an klassische Remote-Meetings denken, in denen eine Entscheidung getroffen
    werden soll: Was läuft dort Ihrer Erfahrung nach häufig schief?
    \begin{itemize}[label=\textendash]
      \item Welche dieser Schwierigkeiten sind für Sie spezifisch für das Remote-Setting und welche würden auch
        in Präsenz auftreten?
      \item \textit{Folgende Stichwörter nicht für explizite Nachfragen nutzen, wenn diese nicht bereits von
        selbst vom Interviewee thematisert wurden, da sonst zu starker Bias \& Risiko nur bereites gewusstes zu finden}
        Zielklarheit und Struktur; Abschweifen vom Thema; ungleiche Beteiligung oder hierarchische Dynamiken;
        Zeitdruck; Ziele, die sich während des Meetings verschieben.
    \end{itemize}
  \item Woran erkennen Sie diese kritischen Momente in denen die gruppe den Faden verliert? Können Sie das
    möglichst an konkreten Beobachtungen festmachen z.B. daran wie gesprochen wird oder wer sprichtoder
    schweigt oder sonstige Zeichen?
    \begin{itemize}[label=\textendash]
      \item \textit{Optional, nicht als Messzwang formulieren:} Gibt es Anzeichen, die
        sich für Sie vergleichsweise zuverlässig wiederholen, sodass Sie fast schon
        sagen könnten: Jetzt muss ich eingreifen?
    \end{itemize}

\end{enumerate}

% ---------------------------------------------------------------------------------------------------------------------
\subsection*{Facilitating Strategien (RQ1)}
\textit{Funktion: Rekonstruktion bewährter Praktiken (Methoden, Strukturen, Gewohnheiten) (Ca. 6 Minuten)}

\begin{enumerate}[label=3.\arabic*, leftmargin=*]
  \item Wie gehen Sie vor, damit ein solches Meeting zielorientiert bleibt und die Gruppe zu einer
    Entscheidung kommt? Beschreiben Sie gerne ein konkretes Vorgehen anhand einer
    typischen Situation.
    \begin{itemize}[label=\textendash]
      \item \textit{Optional:} Welche Methoden, Strukturen oder Gewohnheiten haben sich
        für Sie bewährt (z.\,B.\ Timeboxing, Rollen, Visualisierung, Entscheidungsregeln)?
      \item \textit{Optional:} Was davon legen Sie in der Vorbereitung fest, und was
        steuern Sie erst in der Situation?
    \end{itemize}
  \item Wie stellen Sie sicher, dass den Teilnehmenden während des Meetings klar bleibt,
    was in diesem Moment das eigentliche Ziel ist?
    \begin{itemize}[label=\textendash]
      \item \textit{Optional:} Was geschieht, wenn das Ziel selbst unklar ist oder sich
        im Meeting als unpassend erweist?
    \end{itemize}
  \item Wie wird sicher gestellt dass alle relevanten Informationen die zur erreichung des Ziels wichtig
    sind zur Ansprache kommen? Woher wissen Sie, ob alle relevanten Informationen zur Ansprache kamen?
  \item Wie entscheident ist es, dass alle Teilnehmen zu Wort kommen?
    \begin{itemize}[label=\textendash]
      \item \textit{Falls als relevant erachtet:} Wie stellen Sie sicher das alle Teilnehmen zu Wort kommen?
    \end{itemize}
  \item Worauf achten Sie, wenn sie in das Meeting z.B. in Diskussionen eingreifen? Wie gehen Sie dabei vor?
  \item  Welche Risiken gibt es wenn Sie zu früh oder zu spät eingreifen? Wie versuchen Sie das Risiko zu verringern?
\end{enumerate}

% ---------------------------------------------------------------------------------------------------------------------
\subsection*{Anforderungen an ein KI-gestütztes Moderationssystem (RQ2)}
\textit{Funktion: Offene Anforderungserhebung aus Expertensicht (Ca. 7 Minuten)}

\begin{enumerate}[label=6.\arabic*, leftmargin=*]
  \item Stellen Sie sich ein KI-System vor, das in Remote-Meetings die Teilnehmenden
    dabei unterstützt zielorientiert vorzugehen. Wie sollte ein solches System aus Ihrere Sicht aussehen,
    wenn keine in Facilitating geschulte Person anwesend ist?
  \item Welche Funktionen wären für Sie essenziell, damit das System dabei nützlich ist converget zur
    Entscheidung zu kommen?
  \item Wie sollte dieses System aggieren / intervenieren und was sollte es explizit nicht tun?
  \item Welche Risiken, Schwächen oder unbeabsichtigten Auswirkungen sehen Sie, wenn ein
    KI-System von sich aus moderierend in Meetings einschreitet?
    \begin{itemize}[label=\textendash]
      \item Wie könnte man diese Risiken abschwächen?
      \item \textit{Mögliche Stichwörter:} Vertrauen und Kontrolle der
        Gruppe; Fehlalarme; Gesichtswahrung; hierarchische Verzerrung; Datenschutz und
        Transparenz.
    \end{itemize}
\end{enumerate}

% ---------------------------------------------------------------------------------------------------------------------
\subsection*{Abschluss}

\begin{enumerate}[label=7.\arabic*, leftmargin=*]
  \item Gibt es Aspekte der Zielorientierung, der Entscheidungsfindung oder der
    KI-Unterstützung, die wir nicht angesprochen haben und die Ihnen wichtig erscheinen?
\end{enumerate}

\medskip
